\nivelbasico


\begin{ejercicio}
Una partícula se encuentra en un plano cartesiano en la posición:

\[
(x,y)=(\SI{3.0}{\meter},\SI{4.0}{\meter}).
\]

Determine:

\begin{enumerate}[label=\alph*)]
    \item el vector de posición;
    \item la magnitud del vector de posición;
    \item el ángulo que forma con el eje \(+x\).
\end{enumerate}

\textit{Pista: escriba primero el vector mediante sus componentes rectangulares.}
\end{ejercicio}

\begin{respuesta}
\[
\boxed{\vect{r}=(3.0\,\ihat+4.0\,\jhat)\,\si{\meter}},
\]

\[
\boxed{|\vect{r}|=\SI{5.0}{\meter}},
\]

\[
\boxed{\theta\approx53.1\grado}.
\]
\end{respuesta}

\begin{explicacion}
La posición de una partícula en dos dimensiones se representa mediante:

\[
\vect{r}=x\ihat+y\jhat.
\]

Como:

\[
x=\SI{3.0}{\meter},
\qquad
y=\SI{4.0}{\meter},
\]

tenemos:

\[
\boxed{
\vect{r}
=
(3.0\,\ihat+4.0\,\jhat)\,\si{\meter}
}.
\]

La magnitud del vector de posición se obtiene mediante:

\[
|\vect{r}|
=
\sqrt{x^2+y^2}.
\]

Sustituyendo:

\[
|\vect{r}|
=
\sqrt{(3.0)^2+(4.0)^2}\,\si{\meter}.
\]

Por tanto:

\[
\boxed{|\vect{r}|=\SI{5.0}{\meter}}.
\]

Para determinar la dirección:

\[
\tan\theta=\frac{y}{x}.
\]

Entonces:

\[
\theta
=
\arctan\left(\frac{4.0}{3.0}\right).
\]

Así:

\[
\boxed{\theta\approx53.1\grado}.
\]

La magnitud de \(\vect{r}\) representa la distancia de la partícula al origen en ese instante. No representa la distancia total que haya recorrido para llegar hasta allí.
\end{explicacion}


\begin{ejercicio}
Una partícula se mueve desde el punto:

\[
A(-2,3)\,\si{\meter}
\]

hasta:

\[
B(4,-5)\,\si{\meter}.
\]

Determine:

\begin{enumerate}[label=\alph*)]
    \item los vectores de posición inicial y final;
    \item el vector desplazamiento;
    \item la magnitud del desplazamiento.
\end{enumerate}

\textit{Pista: utilice \(\Delta\vect{r}=\vect{r}_f-\vect{r}_i\) componente a componente.}
\end{ejercicio}

\begin{respuesta}
\[
\vect{r}_i=(-2\,\ihat+3\,\jhat)\,\si{\meter},
\]

\[
\vect{r}_f=(4\,\ihat-5\,\jhat)\,\si{\meter},
\]

\[
\boxed{\Delta\vect{r}=(6\,\ihat-8\,\jhat)\,\si{\meter}},
\]

\[
\boxed{|\Delta\vect{r}|=\SI{10}{\meter}}.
\]
\end{respuesta}

\begin{explicacion}
El vector de posición inicial es:

\[
\vect{r}_i
=
(-2\,\ihat+3\,\jhat)\,\si{\meter}.
\]

El vector de posición final es:

\[
\vect{r}_f
=
(4\,\ihat-5\,\jhat)\,\si{\meter}.
\]

El desplazamiento se define como:

\[
\Delta\vect{r}
=
\vect{r}_f-\vect{r}_i.
\]

Entonces:

\[
\Delta\vect{r}
=
(4-(-2))\ihat
+
(-5-3)\jhat.
\]

Por tanto:

\[
\boxed{
\Delta\vect{r}
=
(6\,\ihat-8\,\jhat)\,\si{\meter}
}.
\]

Su magnitud es:

\[
|\Delta\vect{r}|
=
\sqrt{6^2+(-8)^2}\,\si{\meter}.
\]

Así:

\[
|\Delta\vect{r}|
=
\sqrt{100}\,\si{\meter}.
\]

Por tanto:

\[
\boxed{|\Delta\vect{r}|=\SI{10}{\meter}}.
\]

El desplazamiento depende únicamente de las posiciones inicial y final, no de la trayectoria seguida entre ambos puntos.
\end{explicacion}


\begin{ejercicio}
En el instante:

\[
t_1=\SI{2.0}{\second},
\]

una partícula se encuentra en:

\[
\vect{r}_1=(1\,\ihat-2\,\jhat)\,\si{\meter}.
\]

En:

\[
t_2=\SI{6.0}{\second},
\]

su posición es:

\[
\vect{r}_2=(13\,\ihat+6\,\jhat)\,\si{\meter}.
\]

Determine:

\begin{enumerate}[label=\alph*)]
    \item el desplazamiento;
    \item la velocidad media;
    \item la magnitud de la velocidad media.
\end{enumerate}

\textit{Pista: la velocidad media tiene la misma dirección que el vector desplazamiento.}
\end{ejercicio}

\begin{respuesta}
\[
\Delta\vect{r}
=
(12\,\ihat+8\,\jhat)\,\si{\meter},
\]

\[
\boxed{
\vect{\bar{v}}
=
(3.0\,\ihat+2.0\,\jhat)\,\si{\meter\per\second}
},
\]

\[
\boxed{|\vect{\bar{v}}|\approx\SI{3.61}{\meter\per\second}}.
\]
\end{respuesta}

\begin{explicacion}
El desplazamiento es:

\[
\Delta\vect{r}
=
\vect{r}_2-\vect{r}_1.
\]

Por componentes:

\[
\Delta\vect{r}
=
(13-1)\ihat
+
(6-(-2))\jhat.
\]

Entonces:

\[
\boxed{
\Delta\vect{r}
=
(12\,\ihat+8\,\jhat)\,\si{\meter}
}.
\]

El intervalo de tiempo es:

\[
\Delta t=t_2-t_1.
\]

Así:

\[
\Delta t
=
\SI{6.0}{\second}
-
\SI{2.0}{\second}
=
\SI{4.0}{\second}.
\]

La velocidad media es:

\[
\vect{\bar{v}}
=
\frac{\Delta\vect{r}}{\Delta t}.
\]

Por tanto:

\[
\vect{\bar{v}}
=
\frac{
(12\,\ihat+8\,\jhat)\,\si{\meter}
}
{\SI{4.0}{\second}}.
\]

Entonces:

\[
\boxed{
\vect{\bar{v}}
=
(3.0\,\ihat+2.0\,\jhat)\,\si{\meter\per\second}
}.
\]

Su magnitud es:

\[
|\vect{\bar{v}}|
=
\sqrt{(3.0)^2+(2.0)^2}\,
\si{\meter\per\second}.
\]

Así:

\[
\boxed{
|\vect{\bar{v}}|
=
\sqrt{13}\,\si{\meter\per\second}
\approx
\SI{3.61}{\meter\per\second}
}.
\]
\end{explicacion}


\begin{ejercicio}
En cierto instante, la velocidad de una partícula es:

\[
\vect{v}
=
(6.0\,\ihat-8.0\,\jhat)\,\si{\meter\per\second}.
\]

Determine:

\begin{enumerate}[label=\alph*)]
    \item las componentes de la velocidad;
    \item la rapidez;
    \item el sentido de cada componente.
\end{enumerate}

\textit{Pista: la rapidez es la magnitud del vector velocidad.}
\end{ejercicio}

\begin{respuesta}
\[
v_x=\SI{6.0}{\meter\per\second},
\qquad
v_y=-\SI{8.0}{\meter\per\second},
\]

\[
\boxed{v=\SI{10.0}{\meter\per\second}}.
\]

La componente \(x\) apunta hacia \(+x\) y la componente \(y\) hacia \(-y\).
\end{respuesta}

\begin{explicacion}
Del vector:

\[
\vect{v}
=
(6.0\,\ihat-8.0\,\jhat)\,
\si{\meter\per\second},
\]

identificamos:

\[
v_x=\SI{6.0}{\meter\per\second},
\]

y:

\[
v_y=-\SI{8.0}{\meter\per\second}.
\]

La rapidez es:

\[
v=|\vect{v}|
=
\sqrt{v_x^2+v_y^2}.
\]

Entonces:

\[
v
=
\sqrt{(6.0)^2+(-8.0)^2}
\,\si{\meter\per\second}.
\]

Por tanto:

\[
\boxed{v=\SI{10.0}{\meter\per\second}}.
\]

El signo positivo de \(v_x\) indica movimiento hacia el sentido \(+x\), mientras que el signo negativo de \(v_y\) indica movimiento hacia \(-y\).

La rapidez no lleva signo porque es una magnitud escalar.
\end{explicacion}


\begin{ejercicio}
El vector de posición de una partícula es:

\[
\vect{r}(t)
=
(2+4t)\ihat
+
(3-t^2)\jhat,
\]

donde las posiciones se expresan en metros y el tiempo en segundos.

Determine:

\begin{enumerate}[label=\alph*)]
    \item el vector velocidad como función del tiempo;
    \item el vector aceleración;
    \item la velocidad en \(t=\SI{2.0}{\second}\);
    \item la rapidez en ese instante.
\end{enumerate}

\textit{Pista: derive cada componente del vector de posición respecto del tiempo.}
\end{ejercicio}

\begin{respuesta}
\[
\boxed{
\vect{v}(t)
=
4\,\ihat-2t\,\jhat
}
\]

en \(\si{\meter\per\second}\),

\[
\boxed{
\vect{a}
=
-2\,\jhat\,
\si{\meter\per\second\squared}
},
\]

\[
\vect{v}(\SI{2}{\second})
=
(4\,\ihat-4\,\jhat)\,
\si{\meter\per\second},
\]

\[
\boxed{v\approx\SI{5.66}{\meter\per\second}}.
\]
\end{respuesta}

\begin{explicacion}
La velocidad instantánea es:

\[
\vect{v}
=
\deriv{\vect{r}}{t}.
\]

Derivando cada componente:

\[
v_x
=
\deriv{}{t}(2+4t)
=
4,
\]

\[
v_y
=
\deriv{}{t}(3-t^2)
=
-2t.
\]

Por tanto:

\[
\boxed{
\vect{v}(t)
=
4\,\ihat-2t\,\jhat
}
\]

con unidades de \(\si{\meter\per\second}\).

La aceleración es:

\[
\vect{a}
=
\deriv{\vect{v}}{t}.
\]

Entonces:

\[
a_x=0,
\qquad
a_y=-2.
\]

Así:

\[
\boxed{
\vect{a}
=
-2\,\jhat\,
\si{\meter\per\second\squared}
}.
\]

En:

\[
t=\SI{2.0}{\second},
\]

la velocidad es:

\[
\vect{v}
=
(4\,\ihat-4\,\jhat)\,
\si{\meter\per\second}.
\]

Su magnitud es:

\[
v
=
\sqrt{4^2+(-4)^2}.
\]

Por tanto:

\[
\boxed{
v
=
4\sqrt{2}\,\si{\meter\per\second}
\approx
\SI{5.66}{\meter\per\second}
}.
\]
\end{explicacion}


\begin{ejercicio}
Una partícula tiene inicialmente la velocidad:

\[
\vect{v}_1
=
(3\,\ihat+4\,\jhat)\,
\si{\meter\per\second}.
\]

Dos segundos después su velocidad es:

\[
\vect{v}_2
=
(3\,\ihat-2\,\jhat)\,
\si{\meter\per\second}.
\]

Determine:

\begin{enumerate}[label=\alph*)]
    \item el cambio de velocidad;
    \item la aceleración media;
    \item si la aceleración media es cero porque la componente \(x\) de la velocidad no cambió.
\end{enumerate}

\textit{Pista: la aceleración depende del cambio del vector velocidad completo.}
\end{ejercicio}

\begin{respuesta}
\[
\Delta\vect{v}
=
-6\,\jhat\,
\si{\meter\per\second},
\]

\[
\boxed{
\vect{\bar{a}}
=
-3\,\jhat\,
\si{\meter\per\second\squared}
}.
\]

La aceleración media no es cero.
\end{respuesta}

\begin{explicacion}
El cambio de velocidad es:

\[
\Delta\vect{v}
=
\vect{v}_2-\vect{v}_1.
\]

Entonces:

\[
\Delta\vect{v}
=
(3-3)\ihat
+
(-2-4)\jhat.
\]

Por tanto:

\[
\boxed{
\Delta\vect{v}
=
-6\,\jhat\,
\si{\meter\per\second}
}.
\]

La aceleración media es:

\[
\vect{\bar{a}}
=
\frac{\Delta\vect{v}}{\Delta t}.
\]

Como:

\[
\Delta t=\SI{2}{\second},
\]

tenemos:

\[
\vect{\bar{a}}
=
\frac{-6\,\jhat\,\si{\meter\per\second}}
{\SI{2}{\second}}.
\]

Así:

\[
\boxed{
\vect{\bar{a}}
=
-3\,\jhat\,
\si{\meter\per\second\squared}
}.
\]

Aunque la componente \(x\) de la velocidad permaneció constante, la componente \(y\) cambió.

Por tanto, el vector velocidad cambió y existe aceleración.
\end{explicacion}


\begin{ejercicio}
Una partícula se mueve en tres dimensiones. Su posición cambia desde:

\[
\vect{r}_i
=
(1\,\ihat-2\,\jhat+3\,\khat)\,\si{\meter}
\]

hasta:

\[
\vect{r}_f
=
(5\,\ihat+4\,\jhat-\khat)\,\si{\meter}
\]

en un intervalo de:

\[
\Delta t=\SI{2.0}{\second}.
\]

Determine:

\begin{enumerate}[label=\alph*)]
    \item el desplazamiento;
    \item la magnitud del desplazamiento;
    \item la velocidad media;
    \item la magnitud de la velocidad media.
\end{enumerate}

\textit{Pista: trabaje las tres componentes de forma independiente.}
\end{ejercicio}

\begin{respuesta}
\[
\Delta\vect{r}
=
(4\,\ihat+6\,\jhat-4\,\khat)\,\si{\meter},
\]

\[
|\Delta\vect{r}|
=
\sqrt{68}\,\si{\meter}
\approx
\SI{8.25}{\meter},
\]

\[
\boxed{
\vect{\bar{v}}
=
(2\,\ihat+3\,\jhat-2\,\khat)
\,\si{\meter\per\second}
},
\]

\[
\boxed{
|\vect{\bar{v}}|
=
\sqrt{17}\,\si{\meter\per\second}
\approx
\SI{4.12}{\meter\per\second}
}.
\]
\end{respuesta}

\begin{explicacion}
El desplazamiento es:

\[
\Delta\vect{r}
=
\vect{r}_f-\vect{r}_i.
\]

Por componentes:

\[
\Delta x=5-1=4,
\]

\[
\Delta y=4-(-2)=6,
\]

\[
\Delta z=-1-3=-4.
\]

Por tanto:

\[
\boxed{
\Delta\vect{r}
=
(4\,\ihat+6\,\jhat-4\,\khat)\,\si{\meter}
}.
\]

Su magnitud es:

\[
|\Delta\vect{r}|
=
\sqrt{4^2+6^2+(-4)^2}.
\]

Entonces:

\[
\boxed{
|\Delta\vect{r}|
=
\sqrt{68}\,\si{\meter}
\approx
\SI{8.25}{\meter}
}.
\]

La velocidad media es:

\[
\vect{\bar{v}}
=
\frac{\Delta\vect{r}}{\Delta t}.
\]

Por tanto:

\[
\boxed{
\vect{\bar{v}}
=
(2\,\ihat+3\,\jhat-2\,\khat)
\,\si{\meter\per\second}
}.
\]

Su magnitud:

\[
|\vect{\bar{v}}|
=
\sqrt{2^2+3^2+(-2)^2}.
\]

Así:

\[
\boxed{
|\vect{\bar{v}}|
=
\sqrt{17}\,\si{\meter\per\second}
\approx
\SI{4.12}{\meter\per\second}
}.
\]
\end{explicacion}


\nivelintermedio


\begin{ejercicio}
Una partícula se mueve en el plano \(xy\) con:

\[
\vect{r}_0
=
(2\,\ihat-\jhat)\,\si{\meter},
\]

\[
\vect{v}_0
=
(3\,\ihat+4\,\jhat)\,
\si{\meter\per\second},
\]

y aceleración constante:

\[
\vect{a}
=
(2\,\ihat-\jhat)\,
\si{\meter\per\second\squared}.
\]

Determine después de:

\[
t=\SI{5.0}{\second},
\]

\begin{enumerate}[label=\alph*)]
    \item el vector velocidad;
    \item el vector desplazamiento;
    \item el vector de posición;
    \item la rapidez.
\end{enumerate}

\textit{Pista: aplique las ecuaciones de aceleración constante a cada componente rectangular.}
\end{ejercicio}

\begin{respuesta}
\[
\boxed{
\vect{v}
=
(13\,\ihat-\jhat)\,
\si{\meter\per\second}
},
\]

\[
\boxed{
\Delta\vect{r}
=
(40\,\ihat+7.5\,\jhat)\,\si{\meter}
},
\]

\[
\boxed{
\vect{r}
=
(42\,\ihat+6.5\,\jhat)\,\si{\meter}
},
\]

\[
\boxed{
v=\sqrt{170}\,\si{\meter\per\second}
\approx
\SI{13.0}{\meter\per\second}
}.
\]
\end{respuesta}

\begin{explicacion}
Para aceleración constante:

\[
\vect{v}
=
\vect{v}_0+\vect{a}t.
\]

Sustituyendo:

\[
\vect{v}
=
(3\,\ihat+4\,\jhat)
+
(2\,\ihat-\jhat)(5).
\]

Entonces:

\[
\vect{v}
=
(3+10)\ihat+(4-5)\jhat.
\]

Así:

\[
\boxed{
\vect{v}
=
(13\,\ihat-\jhat)\,
\si{\meter\per\second}
}.
\]

El desplazamiento es:

\[
\Delta\vect{r}
=
\vect{v}_0t+\frac{1}{2}\vect{a}t^2.
\]

La primera contribución es:

\[
\vect{v}_0t
=
(15\,\ihat+20\,\jhat)\,\si{\meter}.
\]

La segunda:

\[
\frac{1}{2}\vect{a}t^2
=
\frac{1}{2}
(2\,\ihat-\jhat)(25).
\]

Entonces:

\[
\frac{1}{2}\vect{a}t^2
=
(25\,\ihat-12.5\,\jhat)\,\si{\meter}.
\]

Sumando:

\[
\boxed{
\Delta\vect{r}
=
(40\,\ihat+7.5\,\jhat)\,\si{\meter}
}.
\]

La posición final es:

\[
\vect{r}
=
\vect{r}_0+\Delta\vect{r}.
\]

Así:

\[
\vect{r}
=
(2+40)\ihat+(-1+7.5)\jhat.
\]

Por tanto:

\[
\boxed{
\vect{r}
=
(42\,\ihat+6.5\,\jhat)\,\si{\meter}
}.
\]

La rapidez es la magnitud de la velocidad:

\[
v
=
\sqrt{13^2+(-1)^2}.
\]

Entonces:

\[
\boxed{
v
=
\sqrt{170}\,\si{\meter\per\second}
\approx
\SI{13.0}{\meter\per\second}
}.
\]
\end{explicacion}


\begin{ejercicio}
Una partícula experimenta aceleración constante. En cierto instante su velocidad es:

\[
\vect{v}_i
=
(2\,\ihat-3\,\jhat)\,
\si{\meter\per\second}.
\]

Cuatro segundos después:

\[
\vect{v}_f
=
(8\,\ihat+5\,\jhat)\,
\si{\meter\per\second}.
\]

Si su posición inicial es:

\[
\vect{r}_i
=
(-5\,\ihat+2\,\jhat)\,\si{\meter},
\]

determine:

\begin{enumerate}[label=\alph*)]
    \item la aceleración;
    \item el desplazamiento durante los cuatro segundos;
    \item la posición final.
\end{enumerate}

\textit{Pista: con aceleración constante, el vector velocidad media es \((\vect{v}_i+\vect{v}_f)/2\).}
\end{ejercicio}

\begin{respuesta}
\[
\boxed{
\vect{a}
=
(1.5\,\ihat+2.0\,\jhat)\,
\si{\meter\per\second\squared}
},
\]

\[
\boxed{
\Delta\vect{r}
=
(20\,\ihat+4\,\jhat)\,\si{\meter}
},
\]

\[
\boxed{
\vect{r}_f
=
(15\,\ihat+6\,\jhat)\,\si{\meter}
}.
\]
\end{respuesta}

\begin{explicacion}
La aceleración constante es:

\[
\vect{a}
=
\frac{\vect{v}_f-\vect{v}_i}{\Delta t}.
\]

El cambio de velocidad es:

\[
\Delta\vect{v}
=
(8-2)\ihat
+
(5-(-3))\jhat.
\]

Entonces:

\[
\Delta\vect{v}
=
(6\,\ihat+8\,\jhat)\,
\si{\meter\per\second}.
\]

Dividiendo entre:

\[
\Delta t=\SI{4}{\second},
\]

obtenemos:

\[
\boxed{
\vect{a}
=
(1.5\,\ihat+2.0\,\jhat)\,
\si{\meter\per\second\squared}
}.
\]

Como la aceleración es constante:

\[
\vect{\bar{v}}
=
\frac{\vect{v}_i+\vect{v}_f}{2}.
\]

Entonces:

\[
\vect{\bar{v}}
=
\frac{
(2+8)\ihat+(-3+5)\jhat
}{2}.
\]

Por tanto:

\[
\vect{\bar{v}}
=
(5\,\ihat+\jhat)\,
\si{\meter\per\second}.
\]

El desplazamiento es:

\[
\Delta\vect{r}
=
\vect{\bar{v}}\Delta t.
\]

Así:

\[
\boxed{
\Delta\vect{r}
=
(20\,\ihat+4\,\jhat)\,\si{\meter}
}.
\]

Finalmente:

\[
\vect{r}_f
=
\vect{r}_i+\Delta\vect{r}.
\]

Por tanto:

\[
\vect{r}_f
=
(-5+20)\ihat
+
(2+4)\jhat.
\]

Así:

\[
\boxed{
\vect{r}_f
=
(15\,\ihat+6\,\jhat)\,\si{\meter}
}.
\]
\end{explicacion}


\begin{ejercicio}
El movimiento de una partícula está descrito por:

\[
\vect{r}(t)
=
2t\,\ihat
+
(t^2-4t+1)\,\jhat,
\]

con posiciones en metros y tiempo en segundos.

Determine:

\begin{enumerate}[label=\alph*)]
    \item los vectores de posición en \(t=\SI{1}{\second}\) y \(t=\SI{3}{\second}\);
    \item la velocidad media durante ese intervalo;
    \item la velocidad instantánea en \(t=\SI{2}{\second}\);
    \item la aceleración;
    \item la ecuación de la trayectoria \(y(x)\).
\end{enumerate}

\textit{Pista: para obtener la trayectoria elimine el tiempo entre \(x(t)\) y \(y(t)\).}
\end{ejercicio}

\begin{respuesta}
\[
\vect{r}(1)
=
(2\,\ihat-2\,\jhat)\,\si{\meter},
\]

\[
\vect{r}(3)
=
(6\,\ihat-2\,\jhat)\,\si{\meter},
\]

\[
\boxed{
\vect{\bar{v}}
=
2\,\ihat\,
\si{\meter\per\second}
},
\]

\[
\boxed{
\vect{v}(2)
=
2\,\ihat\,
\si{\meter\per\second}
},
\]

\[
\boxed{
\vect{a}
=
2\,\jhat\,
\si{\meter\per\second\squared}
},
\]

\[
\boxed{
y=\frac{x^2}{4}-2x+1
}.
\]
\end{respuesta}

\begin{explicacion}
En:

\[
t=\SI{1}{\second},
\]

tenemos:

\[
x=2,
\]

\[
y=1-4+1=-2.
\]

Por tanto:

\[
\vect{r}(1)
=
(2\,\ihat-2\,\jhat)\,\si{\meter}.
\]

En:

\[
t=\SI{3}{\second},
\]

\[
x=6,
\]

\[
y=9-12+1=-2.
\]

Entonces:

\[
\vect{r}(3)
=
(6\,\ihat-2\,\jhat)\,\si{\meter}.
\]

El desplazamiento es:

\[
\Delta\vect{r}
=
4\,\ihat\,\si{\meter}.
\]

El intervalo es:

\[
\Delta t=\SI{2}{\second}.
\]

Por tanto:

\[
\boxed{
\vect{\bar{v}}
=
2\,\ihat\,
\si{\meter\per\second}
}.
\]

La velocidad instantánea se obtiene derivando:

\[
\vect{v}(t)
=
2\,\ihat+(2t-4)\jhat.
\]

En \(t=\SI{2}{\second}\):

\[
\boxed{
\vect{v}(2)
=
2\,\ihat\,
\si{\meter\per\second}
}.
\]

La aceleración es:

\[
\vect{a}
=
\deriv{\vect{v}}{t}
=
2\,\jhat\,
\si{\meter\per\second\squared}.
\]

Para obtener la trayectoria:

\[
x=2t.
\]

Despejamos:

\[
t=\frac{x}{2}.
\]

Sustituimos en:

\[
y=t^2-4t+1.
\]

Entonces:

\[
y
=
\left(\frac{x}{2}\right)^2
-
4\left(\frac{x}{2}\right)
+
1.
\]

Por tanto:

\[
\boxed{
y=\frac{x^2}{4}-2x+1
}.
\]

La trayectoria es una curva, aunque cada componente del movimiento pueda analizarse por separado.
\end{explicacion}


\begin{ejercicio}
Una partícula se mueve de acuerdo con:

\[
x(t)=2t,
\qquad
y(t)=1+t^2,
\]

donde las posiciones están en metros.

Determine:

\begin{enumerate}[label=\alph*)]
    \item la ecuación de la trayectoria;
    \item el vector velocidad;
    \item el vector aceleración;
    \item la posición, velocidad y rapidez en \(t=\SI{3.0}{\second}\).
\end{enumerate}

\textit{Pista: la trayectoria se obtiene eliminando \(t\), mientras que velocidad y aceleración se obtienen derivando respecto del tiempo.}
\end{ejercicio}

\begin{respuesta}
\[
\boxed{
y=1+\frac{x^2}{4}
},
\]

\[
\vect{v}
=
2\,\ihat+2t\,\jhat,
\]

\[
\boxed{
\vect{a}
=
2\,\jhat\,
\si{\meter\per\second\squared}
},
\]

en \(t=\SI{3}{\second}\):

\[
\vect{r}
=
(6\,\ihat+10\,\jhat)\,\si{\meter},
\]

\[
\vect{v}
=
(2\,\ihat+6\,\jhat)\,
\si{\meter\per\second},
\]

\[
\boxed{
v=\sqrt{40}\,\si{\meter\per\second}
\approx
\SI{6.32}{\meter\per\second}
}.
\]
\end{respuesta}

\begin{explicacion}
De:

\[
x=2t,
\]

obtenemos:

\[
t=\frac{x}{2}.
\]

Sustituyendo en:

\[
y=1+t^2,
\]

resulta:

\[
y
=
1+\left(\frac{x}{2}\right)^2.
\]

Por tanto:

\[
\boxed{
y=1+\frac{x^2}{4}
}.
\]

La velocidad tiene componentes:

\[
v_x=\deriv{x}{t}=2,
\]

\[
v_y=\deriv{y}{t}=2t.
\]

Así:

\[
\boxed{
\vect{v}
=
2\,\ihat+2t\,\jhat
}
\]

en \(\si{\meter\per\second}\).

La aceleración es:

\[
a_x=0,
\qquad
a_y=2.
\]

Por tanto:

\[
\boxed{
\vect{a}
=
2\,\jhat\,
\si{\meter\per\second\squared}
}.
\]

En:

\[
t=\SI{3}{\second},
\]

la posición es:

\[
x=2(3)=\SI{6}{\meter},
\]

\[
y=1+(3)^2=\SI{10}{\meter}.
\]

Entonces:

\[
\boxed{
\vect{r}
=
(6\,\ihat+10\,\jhat)\,\si{\meter}
}.
\]

La velocidad:

\[
\vect{v}
=
(2\,\ihat+6\,\jhat)
\,\si{\meter\per\second}.
\]

La rapidez es:

\[
v
=
\sqrt{2^2+6^2}.
\]

Así:

\[
\boxed{
v
=
\sqrt{40}\,\si{\meter\per\second}
\approx
\SI{6.32}{\meter\per\second}
}.
\]
\end{explicacion}


\begin{ejercicio}
Una partícula se mueve en tres dimensiones con vector de posición:

\[
\vect{r}(t)
=
2t\,\ihat
+
t^2\,\jhat
+
(4-t)\,\khat,
\]

donde las posiciones están en metros.

Determine:

\begin{enumerate}[label=\alph*)]
    \item el vector velocidad;
    \item el vector aceleración;
    \item la posición en \(t=\SI{2}{\second}\);
    \item la velocidad y rapidez en ese instante;
    \item la velocidad media entre \(t=\SI{1}{\second}\) y \(t=\SI{3}{\second}\).
\end{enumerate}

\textit{Pista: en tres dimensiones se aplican las mismas definiciones vectoriales componente a componente.}
\end{ejercicio}

\begin{respuesta}
\[
\boxed{
\vect{v}
=
2\,\ihat+2t\,\jhat-\khat
}
\]

en \(\si{\meter\per\second}\),

\[
\boxed{
\vect{a}
=
2\,\jhat\,
\si{\meter\per\second\squared}
},
\]

\[
\vect{r}(2)
=
(4\,\ihat+4\,\jhat+2\,\khat)\,\si{\meter},
\]

\[
\vect{v}(2)
=
(2\,\ihat+4\,\jhat-\khat)\,
\si{\meter\per\second},
\]

\[
\boxed{
v(2)
=
\sqrt{21}\,\si{\meter\per\second}
\approx
\SI{4.58}{\meter\per\second}
},
\]

\[
\boxed{
\vect{\bar{v}}_{1\to3}
=
(2\,\ihat+4\,\jhat-\khat)\,
\si{\meter\per\second}
}.
\]
\end{respuesta}

\begin{explicacion}
Derivamos el vector de posición:

\[
\vect{v}
=
\deriv{\vect{r}}{t}.
\]

Entonces:

\[
\boxed{
\vect{v}
=
2\,\ihat+2t\,\jhat-\khat
}
\]

en \(\si{\meter\per\second}\).

La aceleración:

\[
\vect{a}
=
\deriv{\vect{v}}{t}.
\]

Por tanto:

\[
\boxed{
\vect{a}
=
2\,\jhat\,
\si{\meter\per\second\squared}
}.
\]

En \(t=\SI{2}{\second}\):

\[
\vect{r}(2)
=
4\,\ihat+4\,\jhat+2\,\khat.
\]

Así:

\[
\boxed{
\vect{r}(2)
=
(4\,\ihat+4\,\jhat+2\,\khat)\,\si{\meter}
}.
\]

La velocidad es:

\[
\boxed{
\vect{v}(2)
=
(2\,\ihat+4\,\jhat-\khat)\,
\si{\meter\per\second}
}.
\]

La rapidez:

\[
v
=
\sqrt{2^2+4^2+(-1)^2}.
\]

Por tanto:

\[
\boxed{
v=\sqrt{21}\,\si{\meter\per\second}
\approx
\SI{4.58}{\meter\per\second}
}.
\]

Ahora calculamos las posiciones en \(t=1\) y \(t=3\):

\[
\vect{r}(1)
=
(2\,\ihat+\jhat+3\,\khat)\,\si{\meter},
\]

\[
\vect{r}(3)
=
(6\,\ihat+9\,\jhat+\khat)\,\si{\meter}.
\]

El desplazamiento es:

\[
\Delta\vect{r}
=
(4\,\ihat+8\,\jhat-2\,\khat)\,\si{\meter}.
\]

Como:

\[
\Delta t=\SI{2}{\second},
\]

tenemos:

\[
\boxed{
\vect{\bar{v}}_{1\to3}
=
(2\,\ihat+4\,\jhat-\khat)\,
\si{\meter\per\second}
}.
\]

En este caso la velocidad media entre \(t=1\) y \(t=3\) coincide con la velocidad instantánea en \(t=2\), consecuencia de la aceleración constante.
\end{explicacion}


\begin{ejercicio}
Un vehículo automatizado se mueve sobre un piso horizontal. En \(t=0\):

\[
\vect{r}_0=\vect{0},
\]

y:

\[
\vect{v}_0
=
(4\,\ihat+\jhat)\,
\si{\meter\per\second}.
\]

Su aceleración constante es:

\[
\vect{a}
=
2\,\jhat\,
\si{\meter\per\second\squared}.
\]

Determine después de \(\SI{3.0}{\second}\):

\begin{enumerate}[label=\alph*)]
    \item la posición;
    \item la velocidad;
    \item la rapidez;
    \item la ecuación de la trayectoria \(y(x)\).
\end{enumerate}

\textit{Pista: el movimiento en \(x\) y en \(y\) puede estudiarse de manera independiente, pero ambos tienen el mismo tiempo.}
\end{ejercicio}

\begin{respuesta}
\[
\boxed{
\vect{r}(3)
=
(12\,\ihat+12\,\jhat)\,\si{\meter}
},
\]

\[
\boxed{
\vect{v}(3)
=
(4\,\ihat+7\,\jhat)\,
\si{\meter\per\second}
},
\]

\[
\boxed{
v=\sqrt{65}\,\si{\meter\per\second}
\approx
\SI{8.06}{\meter\per\second}
},
\]

\[
\boxed{
y=\frac{x}{4}+\frac{x^2}{16}
}.
\]
\end{respuesta}

\begin{explicacion}
En la dirección \(x\):

\[
a_x=0.
\]

Por tanto:

\[
x=v_{0x}t=4t.
\]

En la dirección \(y\):

\[
v_{0y}=\SI{1}{\meter\per\second},
\qquad
a_y=\SI{2}{\meter\per\second\squared}.
\]

Entonces:

\[
y
=
v_{0y}t+\frac{1}{2}a_yt^2.
\]

Así:

\[
y=t+t^2.
\]

Para:

\[
t=\SI{3}{\second},
\]

obtenemos:

\[
x=4(3)=\SI{12}{\meter},
\]

\[
y=3+9=\SI{12}{\meter}.
\]

Por tanto:

\[
\boxed{
\vect{r}(3)
=
(12\,\ihat+12\,\jhat)\,\si{\meter}
}.
\]

La componente \(x\) de la velocidad permanece constante:

\[
v_x=\SI{4}{\meter\per\second}.
\]

La componente \(y\):

\[
v_y
=
v_{0y}+a_yt.
\]

Entonces:

\[
v_y
=
1+2(3)
=
\SI{7}{\meter\per\second}.
\]

Por tanto:

\[
\boxed{
\vect{v}(3)
=
(4\,\ihat+7\,\jhat)\,
\si{\meter\per\second}
}.
\]

La rapidez es:

\[
v
=
\sqrt{4^2+7^2}.
\]

Así:

\[
\boxed{
v=\sqrt{65}\,\si{\meter\per\second}
\approx
\SI{8.06}{\meter\per\second}
}.
\]

Para obtener la trayectoria:

\[
x=4t
\quad\Longrightarrow\quad
t=\frac{x}{4}.
\]

Sustituyendo en:

\[
y=t+t^2,
\]

resulta:

\[
y
=
\frac{x}{4}
+
\left(\frac{x}{4}\right)^2.
\]

Por tanto:

\[
\boxed{
y=\frac{x}{4}+\frac{x^2}{16}
}.
\]
\end{explicacion}


\begin{ejercicio}
Una partícula tiene inicialmente:

\[
\vect{v}_i
=
(6\,\ihat+8\,\jhat)\,
\si{\meter\per\second}.
\]

Después de \(\SI{2.0}{\second}\):

\[
\vect{v}_f
=
(-8\,\ihat+6\,\jhat)\,
\si{\meter\per\second}.
\]

Determine:

\begin{enumerate}[label=\alph*)]
    \item la rapidez inicial;
    \item la rapidez final;
    \item el cambio de velocidad;
    \item la aceleración media;
    \item si puede afirmarse que la aceleración es cero porque la rapidez no cambió.
\end{enumerate}

\textit{Pista: velocidad y rapidez no representan la misma magnitud física.}
\end{ejercicio}

\begin{respuesta}
\[
v_i=v_f=\SI{10}{\meter\per\second},
\]

\[
\Delta\vect{v}
=
(-14\,\ihat-2\,\jhat)\,
\si{\meter\per\second},
\]

\[
\boxed{
\vect{\bar{a}}
=
(-7\,\ihat-\jhat)\,
\si{\meter\per\second\squared}
}.
\]

La aceleración media no es cero.
\end{respuesta}

\begin{explicacion}
La rapidez inicial es:

\[
v_i
=
\sqrt{6^2+8^2}.
\]

Entonces:

\[
\boxed{v_i=\SI{10}{\meter\per\second}}.
\]

La rapidez final es:

\[
v_f
=
\sqrt{(-8)^2+6^2}.
\]

Por tanto:

\[
\boxed{v_f=\SI{10}{\meter\per\second}}.
\]

Aunque las magnitudes son iguales, los vectores no lo son.

El cambio de velocidad es:

\[
\Delta\vect{v}
=
\vect{v}_f-\vect{v}_i.
\]

Entonces:

\[
\Delta\vect{v}
=
(-8-6)\ihat
+
(6-8)\jhat.
\]

Así:

\[
\boxed{
\Delta\vect{v}
=
(-14\,\ihat-2\,\jhat)\,
\si{\meter\per\second}
}.
\]

La aceleración media es:

\[
\vect{\bar{a}}
=
\frac{\Delta\vect{v}}{\Delta t}.
\]

Por tanto:

\[
\boxed{
\vect{\bar{a}}
=
(-7\,\ihat-\jhat)\,
\si{\meter\per\second\squared}
}.
\]

Este resultado muestra que una partícula puede tener aceleración aunque su rapidez no cambie.

La aceleración existe porque cambió la dirección del vector velocidad.
\end{explicacion}


\nivelavanzado


\begin{ejercicio}
Una partícula se mueve en el plano con aceleración constante. Se conocen las siguientes posiciones:

\[
\vect{r}(0)
=
(\ihat+2\,\jhat)\,\si{\meter},
\]

\[
\vect{r}(\SI{1}{\second})
=
(5\,\ihat+2\,\jhat)\,\si{\meter},
\]

\[
\vect{r}(\SI{2}{\second})
=
(13\,\ihat+4\,\jhat)\,\si{\meter}.
\]

Determine:

\begin{enumerate}[label=\alph*)]
    \item la velocidad inicial;
    \item la aceleración;
    \item la posición en \(t=\SI{3}{\second}\);
    \item la velocidad en \(t=\SI{3}{\second}\).
\end{enumerate}
\end{ejercicio}

\begin{respuesta}
\[
\boxed{
\vect{v}_0
=
(2\,\ihat-\jhat)\,
\si{\meter\per\second}
},
\]

\[
\boxed{
\vect{a}
=
(4\,\ihat+2\,\jhat)\,
\si{\meter\per\second\squared}
},
\]

\[
\boxed{
\vect{r}(3)
=
(25\,\ihat+8\,\jhat)\,\si{\meter}
},
\]

\[
\boxed{
\vect{v}(3)
=
(14\,\ihat+5\,\jhat)\,
\si{\meter\per\second}
}.
\]
\end{respuesta}

\begin{explicacion}
Para aceleración constante:

\[
\vect{r}(t)
=
\vect{r}_0
+
\vect{v}_0t
+
\frac{1}{2}\vect{a}t^2.
\]

Para:

\[
t=\SI{1}{\second},
\]

tenemos:

\[
\vect{r}(1)-\vect{r}_0
=
\vect{v}_0+\frac{1}{2}\vect{a}.
\]

El desplazamiento es:

\[
(5-1)\ihat+(2-2)\jhat
=
4\,\ihat.
\]

Por tanto:

\[
\vect{v}_0+\frac{1}{2}\vect{a}
=
4\,\ihat.
\]

Para:

\[
t=\SI{2}{\second},
\]

tenemos:

\[
\vect{r}(2)-\vect{r}_0
=
2\vect{v}_0+2\vect{a}.
\]

El desplazamiento es:

\[
(13-1)\ihat+(4-2)\jhat
=
12\,\ihat+2\,\jhat.
\]

Así:

\[
2\vect{v}_0+2\vect{a}
=
12\,\ihat+2\,\jhat.
\]

Multiplicamos la primera ecuación por 2:

\[
2\vect{v}_0+\vect{a}
=
8\,\ihat.
\]

Restando esta ecuación de la segunda:

\[
\vect{a}
=
(12\,\ihat+2\,\jhat)-8\,\ihat.
\]

Por tanto:

\[
\boxed{
\vect{a}
=
(4\,\ihat+2\,\jhat)\,
\si{\meter\per\second\squared}
}.
\]

Sustituyendo en:

\[
\vect{v}_0+\frac{1}{2}\vect{a}
=
4\,\ihat,
\]

obtenemos:

\[
\vect{v}_0
+
(2\,\ihat+\jhat)
=
4\,\ihat.
\]

Entonces:

\[
\boxed{
\vect{v}_0
=
(2\,\ihat-\jhat)\,
\si{\meter\per\second}
}.
\]

Ahora calculamos la posición en \(t=\SI{3}{\second}\):

\[
\vect{r}(3)
=
\vect{r}_0
+
3\vect{v}_0
+
\frac{9}{2}\vect{a}.
\]

Sustituyendo:

\[
\vect{r}(3)
=
(\ihat+2\jhat)
+
(6\ihat-3\jhat)
+
(18\ihat+9\jhat).
\]

Por tanto:

\[
\boxed{
\vect{r}(3)
=
(25\,\ihat+8\,\jhat)\,\si{\meter}
}.
\]

La velocidad es:

\[
\vect{v}(t)
=
\vect{v}_0+\vect{a}t.
\]

Así:

\[
\vect{v}(3)
=
(2\,\ihat-\jhat)
+
3(4\,\ihat+2\,\jhat).
\]

Entonces:

\[
\boxed{
\vect{v}(3)
=
(14\,\ihat+5\,\jhat)\,
\si{\meter\per\second}
}.
\]
\end{explicacion}


\begin{ejercicio}
El movimiento de una partícula está descrito por:

\[
\vect{r}(t)
=
(t^2-4t)\ihat
+
(3t+1)\jhat,
\]

con posiciones en metros.

Determine:

\begin{enumerate}[label=\alph*)]
    \item el vector velocidad;
    \item el vector aceleración;
    \item el instante \(t\geq0\) en que la rapidez es mínima;
    \item la rapidez mínima;
    \item la posición en ese instante;
    \item la relación geométrica entre \(\vect{v}\) y \(\vect{a}\) en el instante de rapidez mínima.
\end{enumerate}
\end{ejercicio}

\begin{respuesta}
\[
\vect{v}
=
(2t-4)\ihat+3\jhat,
\]

\[
\vect{a}
=
2\ihat\,
\si{\meter\per\second\squared},
\]

\[
\boxed{t=\SI{2.0}{\second}},
\]

\[
\boxed{v_{\min}=\SI{3.0}{\meter\per\second}},
\]

\[
\boxed{
\vect{r}(2)
=
(-4\,\ihat+7\,\jhat)\,\si{\meter}
}.
\]

En ese instante:

\[
\boxed{\vect{v}\perp\vect{a}}.
\]
\end{respuesta}

\begin{explicacion}
La velocidad se obtiene derivando:

\[
\vect{v}
=
\deriv{\vect{r}}{t}.
\]

Entonces:

\[
\boxed{
\vect{v}
=
(2t-4)\ihat+3\jhat
}
\]

en \(\si{\meter\per\second}\).

La aceleración es:

\[
\vect{a}
=
\deriv{\vect{v}}{t}.
\]

Por tanto:

\[
\boxed{
\vect{a}
=
2\ihat\,
\si{\meter\per\second\squared}
}.
\]

La rapidez satisface:

\[
v
=
\sqrt{(2t-4)^2+3^2}.
\]

Para encontrar su mínimo es suficiente minimizar:

\[
v^2=(2t-4)^2+9.
\]

El primer término nunca puede ser negativo y alcanza su valor mínimo cuando:

\[
2t-4=0.
\]

Por tanto:

\[
\boxed{t=\SI{2.0}{\second}}.
\]

En ese instante:

\[
\vect{v}(2)
=
3\jhat\,
\si{\meter\per\second}.
\]

Por tanto:

\[
\boxed{v_{\min}=\SI{3.0}{\meter\per\second}}.
\]

La posición es:

\[
x(2)
=
(2)^2-4(2)
=
-4,
\]

\[
y(2)
=
3(2)+1
=
7.
\]

Así:

\[
\boxed{
\vect{r}(2)
=
(-4\,\ihat+7\,\jhat)\,\si{\meter}
}.
\]

Además:

\[
\vect{v}(2)
=
3\jhat,
\]

mientras que:

\[
\vect{a}=2\ihat.
\]

Por tanto:

\[
\vect{v}\pdot\vect{a}=0.
\]

Así:

\[
\boxed{\vect{v}\perp\vect{a}}.
\]

En ese instante la aceleración cambia la dirección del vector velocidad sin producir, instantáneamente, un aumento o disminución de su magnitud.
\end{explicacion}


\begin{ejercicio}
Una partícula se mueve en tres dimensiones con aceleración constante. En \(t=0\):

\[
\vect{r}_0
=
(\jhat+2\,\khat)\,\si{\meter},
\]

\[
\vect{v}_0
=
(\ihat+2\,\jhat-\khat)\,
\si{\meter\per\second},
\]

y:

\[
\vect{a}
=
(2\,\ihat+2\,\khat)\,
\si{\meter\per\second\squared}.
\]

Determine el primer instante \(t\geq0\) en que la partícula alcanza el plano:

\[
x+y+z=\SI{15}{\meter}.
\]

En ese instante determine también:

\begin{enumerate}[label=\alph*)]
    \item el vector de posición;
    \item el vector velocidad;
    \item la rapidez.
\end{enumerate}
\end{ejercicio}

\begin{respuesta}
\[
\boxed{t=\SI{2.0}{\second}},
\]

\[
\boxed{
\vect{r}
=
(6\,\ihat+5\,\jhat+4\,\khat)\,\si{\meter}
},
\]

\[
\boxed{
\vect{v}
=
(5\,\ihat+2\,\jhat+3\,\khat)\,
\si{\meter\per\second}
},
\]

\[
\boxed{
v=\sqrt{38}\,\si{\meter\per\second}
\approx
\SI{6.16}{\meter\per\second}
}.
\]
\end{respuesta}

\begin{explicacion}
Para aceleración constante:

\[
\vect{r}
=
\vect{r}_0
+
\vect{v}_0t
+
\frac{1}{2}\vect{a}t^2.
\]

Por componentes:

\[
x=t+t^2,
\]

\[
y=1+2t,
\]

\[
z=2-t+t^2.
\]

La partícula alcanza el plano cuando:

\[
x+y+z=15.
\]

Sustituyendo:

\[
(t+t^2)+(1+2t)+(2-t+t^2)=15.
\]

Simplificando:

\[
2t^2+2t+3=15.
\]

Entonces:

\[
2t^2+2t-12=0.
\]

Dividiendo entre 2:

\[
t^2+t-6=0.
\]

Factorizando:

\[
(t+3)(t-2)=0.
\]

Las soluciones son:

\[
t=-\SI{3}{\second},
\]

y:

\[
t=\SI{2}{\second}.
\]

Como buscamos:

\[
t\geq0,
\]

la solución física es:

\[
\boxed{t=\SI{2.0}{\second}}.
\]

Evaluamos la posición:

\[
x=2+4=6,
\]

\[
y=1+4=5,
\]

\[
z=2-2+4=4.
\]

Por tanto:

\[
\boxed{
\vect{r}
=
(6\,\ihat+5\,\jhat+4\,\khat)\,\si{\meter}
}.
\]

La velocidad es:

\[
\vect{v}
=
\vect{v}_0+\vect{a}t.
\]

Entonces:

\[
\vect{v}
=
(\ihat+2\jhat-\khat)
+
2(2\ihat+2\khat).
\]

Así:

\[
\boxed{
\vect{v}
=
(5\,\ihat+2\,\jhat+3\,\khat)\,
\si{\meter\per\second}
}.
\]

La rapidez es:

\[
v
=
\sqrt{5^2+2^2+3^2}.
\]

Por tanto:

\[
\boxed{
v
=
\sqrt{38}\,\si{\meter\per\second}
\approx
\SI{6.16}{\meter\per\second}
}.
\]

La raíz negativa es matemáticamente válida para la función de movimiento extendida hacia tiempos anteriores, pero no pertenece al intervalo físico solicitado.
\end{explicacion}


\begin{ejercicio}
Una partícula cambia su velocidad desde:

\[
\vect{v}_i
=
(6\,\ihat+8\,\jhat)\,
\si{\meter\per\second}
\]

hasta:

\[
\vect{v}_f
=
(-8\,\ihat+6\,\jhat)\,
\si{\meter\per\second}
\]

en \(\SI{2.0}{\second}\).

Un estudiante afirma:

\begin{quote}
``Como la rapidez inicial y la rapidez final son iguales, la aceleración media es cero''.
\end{quote}

Analice la afirmación mediante componentes rectangulares y explique con precisión el error conceptual.
\end{ejercicio}

\begin{respuesta}
La afirmación es incorrecta.

Aunque:

\[
|\vect{v}_i|
=
|\vect{v}_f|
=
\SI{10}{\meter\per\second},
\]

el vector velocidad cambió.

Por tanto:

\[
\boxed{
\vect{\bar{a}}
=
(-7\,\ihat-\jhat)\,
\si{\meter\per\second\squared}
}.
\]
\end{respuesta}

\begin{explicacion}
Calculamos primero las magnitudes.

Para la velocidad inicial:

\[
|\vect{v}_i|
=
\sqrt{6^2+8^2}
=
\SI{10}{\meter\per\second}.
\]

Para la velocidad final:

\[
|\vect{v}_f|
=
\sqrt{(-8)^2+6^2}
=
\SI{10}{\meter\per\second}.
\]

La rapidez no cambió.

Sin embargo, la aceleración no se define mediante el cambio de rapidez, sino mediante el cambio del \textbf{vector velocidad}:

\[
\vect{\bar{a}}
=
\frac{\vect{v}_f-\vect{v}_i}{\Delta t}.
\]

El cambio de velocidad es:

\[
\Delta\vect{v}
=
(-8-6)\ihat
+
(6-8)\jhat.
\]

Entonces:

\[
\Delta\vect{v}
=
(-14\,\ihat-2\,\jhat)\,
\si{\meter\per\second}.
\]

Dividiendo entre:

\[
\Delta t=\SI{2.0}{\second},
\]

obtenemos:

\[
\boxed{
\vect{\bar{a}}
=
(-7\,\ihat-\jhat)\,
\si{\meter\per\second\squared}
}.
\]

El error conceptual consiste en confundir \textbf{rapidez} y \textbf{velocidad}.

La rapidez sólo describe la magnitud de \(\vect{v}\). La velocidad incluye además dirección y sentido.

Por tanto, una partícula puede acelerar aunque su rapidez permanezca constante si cambia la dirección de su movimiento.
\end{explicacion}


\begin{ejercicio}
Partiendo de las definiciones:

\[
\vect{v}
=
\deriv{\vect{r}}{t}
\]

y:

\[
\vect{a}
=
\deriv{\vect{v}}{t},
\]

demuestre que, en coordenadas cartesianas con ejes fijos,

\[
\vect{v}
=
v_x\ihat+v_y\jhat+v_z\khat
\]

implica:

\[
v_x=\deriv{x}{t},
\qquad
v_y=\deriv{y}{t},
\qquad
v_z=\deriv{z}{t},
\]

y que:

\[
a_x=\deriv{v_x}{t},
\qquad
a_y=\deriv{v_y}{t},
\qquad
a_z=\deriv{v_z}{t}.
\]

Explique qué consecuencia práctica tiene este resultado para resolver problemas de movimiento en dos y tres dimensiones.
\end{ejercicio}

\begin{respuesta}
Se obtiene:

\[
\boxed{
\vect{v}
=
\deriv{x}{t}\ihat
+
\deriv{y}{t}\jhat
+
\deriv{z}{t}\khat
}
\]

y:

\[
\boxed{
\vect{a}
=
\deriv{v_x}{t}\ihat
+
\deriv{v_y}{t}\jhat
+
\deriv{v_z}{t}\khat
}.
\]

Por tanto, el movimiento puede analizarse componente a componente sobre ejes cartesianos fijos.
\end{respuesta}

\begin{explicacion}
El vector de posición en tres dimensiones es:

\[
\vect{r}
=
x\ihat+y\jhat+z\khat.
\]

La velocidad se define como:

\[
\vect{v}
=
\deriv{\vect{r}}{t}.
\]

Por tanto:

\[
\vect{v}
=
\deriv{}{t}
\left(
x\ihat+y\jhat+z\khat
\right).
\]

En un sistema cartesiano con ejes fijos, los vectores unitarios:

\[
\ihat,\qquad
\jhat,\qquad
\khat
\]

no cambian con el tiempo.

Por tanto:

\[
\deriv{\ihat}{t}
=
\deriv{\jhat}{t}
=
\deriv{\khat}{t}
=
0.
\]

Entonces:

\[
\vect{v}
=
\deriv{x}{t}\ihat
+
\deriv{y}{t}\jhat
+
\deriv{z}{t}\khat.
\]

Comparando con:

\[
\vect{v}
=
v_x\ihat+v_y\jhat+v_z\khat,
\]

obtenemos:

\[
\boxed{
v_x=\deriv{x}{t}
},
\]

\[
\boxed{
v_y=\deriv{y}{t}
},
\]

\[
\boxed{
v_z=\deriv{z}{t}
}.
\]

Ahora:

\[
\vect{a}
=
\deriv{\vect{v}}{t}.
\]

Como:

\[
\vect{v}
=
v_x\ihat+v_y\jhat+v_z\khat,
\]

tenemos:

\[
\vect{a}
=
\deriv{v_x}{t}\ihat
+
\deriv{v_y}{t}\jhat
+
\deriv{v_z}{t}\khat.
\]

Por tanto:

\[
\boxed{
a_x=\deriv{v_x}{t}
},
\]

\[
\boxed{
a_y=\deriv{v_y}{t}
},
\]

\[
\boxed{
a_z=\deriv{v_z}{t}
}.
\]

La consecuencia práctica es fundamental: en coordenadas cartesianas con ejes fijos, un problema vectorial puede descomponerse en problemas escalares independientes para cada componente.

Así, podemos estudiar por separado:

\[
x(t),\qquad y(t),\qquad z(t),
\]

manteniendo siempre un mismo parámetro temporal \(t\).

Después de resolver las componentes, se reconstruyen los vectores completos de posición, velocidad y aceleración.
\end{explicacion}


\begin{ejercicio}
Una partícula tiene velocidad:

\[
\vect{v}(t)
=
(4-2t)\ihat
+
(1+3t)\jhat,
\]

en \(\si{\meter\per\second}\).

En \(t=0\) se encuentra en:

\[
\vect{r}_0
=
(-\ihat+2\,\jhat)\,\si{\meter}.
\]

Determine:

\begin{enumerate}[label=\alph*)]
    \item el vector aceleración;
    \item el vector de posición \(\vect{r}(t)\);
    \item el primer instante \(t\geq0\) en que la velocidad es paralela al vector \(\ihat+\jhat\);
    \item la posición y la rapidez en ese instante.
\end{enumerate}
\end{ejercicio}

\begin{respuesta}
\[
\boxed{
\vect{a}
=
(-2\,\ihat+3\,\jhat)\,
\si{\meter\per\second\squared}
},
\]

\[
\boxed{
\vect{r}(t)
=
(-1+4t-t^2)\ihat
+
\left(2+t+\frac{3}{2}t^2\right)\jhat
}
\]

en metros,

\[
\boxed{t=\SI{0.60}{\second}},
\]

\[
\boxed{
\vect{r}
=
(1.04\,\ihat+3.14\,\jhat)\,\si{\meter}
},
\]

\[
\boxed{
v\approx\SI{3.96}{\meter\per\second}
}.
\]
\end{respuesta}

\begin{explicacion}
La aceleración es:

\[
\vect{a}
=
\deriv{\vect{v}}{t}.
\]

Derivando:

\[
\boxed{
\vect{a}
=
(-2\,\ihat+3\,\jhat)\,
\si{\meter\per\second\squared}
}.
\]

Para obtener la posición integramos cada componente de la velocidad.

En \(x\):

\[
v_x=4-2t.
\]

Entonces:

\[
x(t)
=
x_0+\int_0^t(4-2t')\,\dd t'.
\]

Como:

\[
x_0=-1,
\]

resulta:

\[
x(t)
=
-1+4t-t^2.
\]

En \(y\):

\[
v_y=1+3t.
\]

Entonces:

\[
y(t)
=
y_0+\int_0^t(1+3t')\,\dd t'.
\]

Como:

\[
y_0=2,
\]

obtenemos:

\[
y(t)
=
2+t+\frac{3}{2}t^2.
\]

Por tanto:

\[
\boxed{
\vect{r}(t)
=
(-1+4t-t^2)\ihat
+
\left(2+t+\frac{3}{2}t^2\right)\jhat
}.
\]

Para que \(\vect{v}\) sea paralela a:

\[
\ihat+\jhat,
\]

sus componentes deben ser iguales:

\[
v_x=v_y.
\]

Entonces:

\[
4-2t=1+3t.
\]

Por tanto:

\[
3=5t,
\]

y:

\[
\boxed{t=\SI{0.60}{\second}}.
\]

En ese instante:

\[
v_x
=
4-2(0.60)
=
\SI{2.8}{\meter\per\second},
\]

\[
v_y
=
1+3(0.60)
=
\SI{2.8}{\meter\per\second}.
\]

La rapidez es:

\[
v
=
\sqrt{(2.8)^2+(2.8)^2}.
\]

Así:

\[
\boxed{
v\approx\SI{3.96}{\meter\per\second}
}.
\]

La posición en ese instante es:

\[
x
=
-1+4(0.60)-(0.60)^2.
\]

Entonces:

\[
x
=
-1+2.40-0.36
=
\SI{1.04}{\meter}.
\]

Para \(y\):

\[
y
=
2+0.60+\frac{3}{2}(0.60)^2.
\]

Por tanto:

\[
y
=
2+0.60+0.54
=
\SI{3.14}{\meter}.
\]

Así:

\[
\boxed{
\vect{r}
=
(1.04\,\ihat+3.14\,\jhat)\,\si{\meter}
}.
\]
\end{explicacion}